\documentclass[UTF8,12pt,a4paper]{ctexart}

\usepackage{amsmath,amssymb,booktabs,geometry,enumitem,hyperref}
\geometry{left=2.6cm,right=2.6cm,top=2.7cm,bottom=2.7cm}
\hypersetup{colorlinks=true,linkcolor=black,citecolor=black,urlcolor=black}
\setlist[itemize]{leftmargin=2em,itemsep=0.2em,topsep=0.2em}
\setlist[enumerate]{leftmargin=2em,itemsep=0.2em,topsep=0.2em}

\title{\textbf{JIE 2019 资本积累与动态贸易收益复刻笔记}\\
\large 从 MATLAB 到 Python 的模型求解、校准核验与绘图复现}
\author{复刻整理}
\date{\today}

\begin{document}
\maketitle

\begin{abstract}
这份笔记总结我复刻 Ravikumar、Santacreu 和 Sposi（2019）代码的过程。复刻内容包括数据处理、参数校准、稳态求解、动态路径求解思路、MATLAB 结果核验和 Python 绘图复现。主要结论是：自我复刻的 CSV 校准结果已经与作者 MATLAB 参数文件严格一致；此前所谓“归一化没做好”的问题，主要来自把 MATLAB 中的收入指标 $y=w/[(1-\alpha)P_c]$ 误当成中间品总产出 $Y_m$，而不是原始数据归一化错误。完成这次复刻后，我对大型结构模型代码的理解从“看公式”推进到“能拆模块、能查维度、能核验结果、能画出对应图形”。
\end{abstract}

\tableofcontents
\newpage

\section{复刻目标}

这次复刻围绕 JIE 2019 文章的基准模型展开，核心目标有四个。

\begin{enumerate}
  \item 用 Python 复刻作者 MATLAB 的数据处理、校准和稳态求解。
  \item 检查 Python 结果与作者 MATLAB 结果的归一化和变量定义是否对应。
  \item 用 Python 的 seaborn 和 matplotlib 复现作者 MATLAB 代码中的主要图形。
  \item 总结复刻大型动态结构模型代码的经验，尤其是如何从理论模型走到可运行代码。
\end{enumerate}

我没有把两国模型作为主要复刻对象。两国模型更像是作者用于说明机制或补充验证的低维版本，对我当前最重要的训练价值不如 44 国基准模型、校准体系、过渡动态和图形复现。本文档仍然承认两国模型存在，但不把它作为能力展示的核心。

\section{校准与归一化核验}

\subsection{数据归一化本身没有错}

我用新增脚本 \texttt{jie2019\_check\_and\_plot.py} 读取自我复刻生成的三张 CSV：
\begin{itemize}
  \item \texttt{01\_baseline\_country\_table.csv}
  \item \texttt{02\_baseline\_trade\_table.csv}
  \item \texttt{03\_baseline\_parameter\_table.csv}
\end{itemize}
并逐项比较作者 MATLAB 文件 \texttt{ParametersSimple.mat} 中的对应变量。校验结果输出为：
\[
\texttt{Python\_复刻输出/calibration\_check.csv}.
\]

校验表显示，Python 校准结果与 MATLAB 参数结果在 GDP、消费、投资、部门产出、增加值、贸易份额、贸易成本、国家参数和技术参数上高度一致。最大相对误差约为
\[
2.99\times 10^{-9},
\]
可以视为数值舍入误差。这说明原始数据处理和校准阶段的归一化没有系统性错误。

进一步地，我还生成了稳态对齐表：
\[
\texttt{Python\_复刻输出/steady\_state\_alignment.csv}.
\]
该表比较 Python 稳态求解结果与 MATLAB 初始稳态变量，包括工资、租金、价格、资本、消费、投资、收入指标、TFP、相对价格和贸易份额。核心变量相关系数接近 1，最大相对误差约为 $2.42\times 10^{-4}$，主要来自求解容忍度和迭代方法差异。

\subsection{真正需要修正的是变量含义}

最初我怀疑 Python 的归一化没有做好，是因为把 Python 里的 $Y_m$ 与 MATLAB 汇总结果里的 \texttt{y1\_CfUniformLib} 做了比较，结果差异很大。进一步检查 MATLAB 的 \texttt{MergeBaselineResults.m} 后发现：
\begin{equation}
y_1=\frac{w_1}{(1-\alpha_1)P_{c1}}.
\end{equation}
也就是说，MATLAB 的 \texttt{y1} 是按消费品价格折算的收入或产出指标，不是中间品总产出 $Y_m$。

Python 的中间品产出方程与 MATLAB 的 \texttt{prodn\_alloc.m} 实际一致。中间品产出由线性系统给出：
\begin{equation}
P_{m,i}Y_{m,i}
-
\sum_j \pi_{j i}(1-\nu_{m,j})P_{m,j}Y_{m,j}
=
\sum_j \pi_{j i}\left[(1-\nu_{c,j})P_{c,j}C_j+(1-\nu_{x,j})P_{x,j}X_j\right].
\end{equation}
用 MATLAB 的价格、贸易份额和支出变量代入 Python 公式，可以复现 MATLAB 的贸易赤字 $F$。因此，问题不在 $Y_m$ 的线性方程，而在比较对象选错。

\section{稳态求解理解}

作者的稳态求解不是一次性求所有内生变量，而是把问题压缩到工资向量。给定工资后，可以依次求出：
\begin{enumerate}
  \item 租金、价格指数和贸易份额；
  \item 资本、投资和消费；
  \item 中间品产出和部门要素分配；
  \item 贸易赤字或外部清算残差。
\end{enumerate}

工资迭代的核心残差是外部清算条件。MATLAB 的 \texttt{SS\_compeqbm.m} 用如下方式更新工资：
\begin{equation}
\tilde{w}_i=w_i\left(1+s\frac{Z_i}{L_i}\right),
\end{equation}
其中 $Z_i$ 是外部清算残差。除以 $L_i$ 的作用是保持总劳动收入归一化：
\begin{equation}
\sum_i w_iL_i=\text{常数}.
\end{equation}
Python 版本使用 \texttt{least\_squares} 直接求解工资残差，并通过
\[
w_i=\frac{\exp(\omega_i)}{\sum_j \exp(\omega_j)L_j}
\]
保证同一个计价物。与 MATLAB 初始稳态比较后，工资、价格、租金、消费、投资和资本的相对误差都在很小范围内，说明稳态求解框架与作者代码对应。

\section{动态路径与数值优化}

动态路径求解比稳态困难得多。稳态中许多变量可以顺着均衡条件依次推出；动态模型中，消费、投资、资本、净国外资产和世界利率相互依赖，不能简单逐期求解。

我对动态路径的理解是：这是一个大型非线性边值问题。初始资本和初始资产头寸由基准稳态给定，终端条件由反事实稳态给定，中间每一期都必须满足欧拉方程、资本积累、资产积累、价格指数、贸易份额和市场出清。抽象地写，就是求
\begin{equation}
\bm R(\bm x)=0,
\end{equation}
其中 $\bm x$ 包含所有时期的工资、消费、投资、价格、资本、资产和世界利率。

复刻过程中最大的能力提升来自稀疏结构意识。若把每期变量平铺成一个大向量，许多残差只依赖相邻时期变量，雅可比矩阵天然稀疏。利用 \texttt{jac\_sparsity} 可以减少无效数值差分和线性子问题计算量。这让我理解到，写动态模型代码不是把公式机械翻译成循环，而是要把模型的局部依赖结构告诉求解器。

\section{Python 绘图复现}

我新增了 Python 脚本：
\[
\texttt{jie2019\_check\_and\_plot.py}.
\]
脚本读取 MATLAB 汇总结果 \texttt{MergedBaselineResults.mat} 和 \texttt{MergedStaticXinCResults.mat}，用 seaborn 和 matplotlib 复现作者 \texttt{PlotResults.m} 中的主要图形。输出目录为：
\[
\texttt{Python\_复刻输出/figures}.
\]

目前已生成 20 张图，包括：
\begin{itemize}
  \item tradables intensity 分布；
  \item 20\% 贸易成本下降下的动态收益与 GDP；
  \item 消费、经常账户、TFP、相对价格、收入和资本的过渡路径；
  \item 半衰期与经常账户；
  \item 动态收益弹性；
  \item 统一贸易自由化与单边贸易自由化收益比较；
  \item 初始 NFA 与动态收益相对即时收益；
  \item $\nu_c$ 与 $\nu_x$ 替换实验；
  \item 静态模型与动态模型的 $\nu_c$ 对比；
  \item 投资率分布与开放度图。
\end{itemize}

图形清单写入：
\[
\texttt{Python\_复刻输出/python\_figure\_manifest.csv}.
\]
这些图不追求与 MATLAB 的线型、字号、标签位置完全一致，但数值来源和变量定义与 MATLAB 对应。

\section{代码能力提升}

这次复刻使我对“会写模型代码”有了更具体的理解。

\subsection{从公式到数组}

贸易、空间和多国模型最容易出错的地方不是公式，而是数组维度。一个贸易份额 $\pi_{ij}$ 到底表示“目的国 $i$ 从来源国 $j$ 进口”，还是“来源国 $i$ 出口到目的国 $j$”，会直接决定矩阵是否需要转置。复刻中我学到的经验是：每个数组都要带着经济含义命名，并在写矩阵乘法前先写清楚自由下标和求和下标。

\subsection{从结果到核验}

大型模型不能只看代码是否能运行。必须建立核验表，逐项检查 Python 与 MATLAB 的数据、参数、份额、价格和稳态变量是否对应。校验表比肉眼看图更可靠，因为它能直接指出误差来自数据处理、参数校准、稳态求解还是图形变量定义。

\subsection{从逐行翻译到模块化}

MATLAB 代码常把经济逻辑拆成 \texttt{priceindex}、\texttt{tradeshare}、\texttt{alloc}、\texttt{compeqbm} 等模块。Python 复刻时，我也逐渐理解到，不能把所有内容写进一个大函数。更好的结构是：
\begin{enumerate}
  \item 参数读取与校准；
  \item 价格指数；
  \item 贸易份额；
  \item 数量分配；
  \item 稳态残差；
  \item 动态路径残差；
  \item 结果核验与绘图。
\end{enumerate}
这种模块化使代码更容易定位错误，也更容易替换求解器或加入新冲击。

\subsection{从“调包”到理解求解器}

求解器不是黑箱。残差告诉求解器当前错在哪里，雅可比告诉求解器往哪个方向调整能减少错误。动态模型的求解速度和稳定性，很大程度取决于是否正确缩放变量、是否使用对数变换保证正值、是否利用稀疏雅可比、是否给出合理初值。

\section{方法论总结}

掌握这套复刻方法后，可以处理很多现实问题。只要能把现实问题翻译成状态变量、控制变量、约束和冲击，就可以用类似框架研究：
\begin{itemize}
  \item 贸易成本下降对不同国家资本积累和福利的动态影响；
  \item 产业政策如何通过投资和生产率改变长期产出；
  \item 汇率或外部融资冲击如何通过净国外资产传播；
  \item 基础设施建设如何影响地区间贸易成本和资本配置；
  \item 技术扩散或知识流动如何改变长期增长路径。
\end{itemize}

我现在更能理解老师所说的“模型好做，但代码不好写”。理论模型写在纸上时，很多下标、归一化和均衡条件可以靠直觉略过；真正写代码时，每一个下标方向、每一个价格单位、每一个残差尺度都必须明确。完成这次复刻后，我的提升不只是会运行一段 Python，而是知道如何把结构模型拆成可核验、可调试、可扩展的数值程序。

\section{文件清单}

\begin{itemize}
  \item 校验与绘图脚本：\texttt{jie2019\_check\_and\_plot.py}
  \item 校验表：\texttt{Python\_复刻输出/calibration\_check.csv}
  \item 稳态对齐表：\texttt{Python\_复刻输出/steady\_state\_alignment.csv}
  \item 图形清单：\texttt{Python\_复刻输出/python\_figure\_manifest.csv}
  \item Python 图形目录：\texttt{Python\_复刻输出/figures}
  \item 本文档：\texttt{jie2019\_replication\_notes.tex}
\end{itemize}

\begin{thebibliography}{9}
\bibitem{ravikumar2019}
Ravikumar, B., Santacreu, A. M., and Sposi, M. 2019.
\textit{Capital Accumulation and Dynamic Gains from Trade}.
Journal of International Economics, 119, 93--110.
\end{thebibliography}

\end{document}
